\documentclass[conference]{IEEEtran}
\IEEEoverridecommandlockouts

\usepackage{cite}
\usepackage{amsmath,amssymb}
\usepackage{graphicx}
\usepackage{xcolor}
\usepackage{url}
% Local: figures live in runs/workshop2/figures/  Overleaf: figures/ (see scripts/pack_report.py)
\graphicspath{{../runs/workshop2/figures/}{figures/}}

% ---------------------------------------------------------------------------
% Helper macros
% \todo{...}  -> red text for things the authors must fill in
% \figph[width]{description}{file} -> inserts figures/<file>.png when it exists
%               (generate them with: python src/w2/make_figures_w2.py); otherwise
%               it shows a framed placeholder so the document always compiles.
% ---------------------------------------------------------------------------
\newcommand{\todo}[1]{\textcolor{red}{\textbf{[TODO: #1]}}}
\newcommand{\figph}[3][\columnwidth]{%
  \includegraphics[width=#1]{#3.png}}

\begin{document}

\title{Deep Sequence Model for Failure Prediction in CartPole Trajectories}

\author{\IEEEauthorblockN{Miguel Angel Torres Ramirez}
\IEEEauthorblockA{Dept. of Computer Engineering\\
Universidad Distrital Francisco Jos\'e de Caldas\\
Email: migatorresr@udistrital.edu.co}
\and
\IEEEauthorblockN{Cristian David Romero Gil}
\IEEEauthorblockA{Dept. of Computer Engineering\\
Universidad Distrital Francisco Jos\'e de Caldas\\
Email: cdromerog@udistrital.edu.co}
\and
\IEEEauthorblockN{Omar Esteban Peña Sanchez}
\IEEEauthorblockA{Dept. of Computer Engineering\\
Universidad Distrital Francisco Jos\'e de Caldas\\
Email: oepenas@udistrital.edu.co}}

\maketitle

\begin{abstract}
This report presents the deep learning model developed in Workshop 2 for the
binary failure-prediction task on CartPole trajectories (35,874 samples from
400 episodes, 17.2:1 class imbalance). Because the label of a row depends on
how the system evolves over the next few steps, we selected a recurrent
architecture: a 2-layer LSTM (64 hidden units, 53{,}569 parameters) that reads
a past-only window of the last 16 steps of each episode. Training combines
position-weighted binary cross-entropy for the imbalance, AdamW with a cosine
learning-rate schedule, dropout, weight decay, gradient clipping and early
stopping on validation F1; the decision threshold is selected on the
validation set, exactly as in Workshop 1, so the test set is never used for
tuning. On the held-out test episodes the model reaches F1 = 0.865,
PR-AUC = 0.956 and recall = 0.924, improving on the best Workshop 1 baseline
(Gradient Boosting: F1 = 0.860, PR-AUC = 0.935, recall = 0.886) while using
only the five raw state variables -- the network learns its own temporal
features instead of relying on the hand-crafted ones.
\end{abstract}

\begin{IEEEkeywords}
CartPole, failure prediction, LSTM, deep learning, class imbalance, sequence classification.
\end{IEEEkeywords}

%% ======================================================================
\section{Introduction}
Workshop 1 established the difficulty of the failure-prediction problem on
CartPole trajectories and the reference performance of three traditional
classifiers (F1 up to 0.860 for Gradient Boosting) using hand-crafted
features~\cite{barto1983}. Those models, however, classify each row
\emph{independently}: they see a single state and must infer the dynamics
from engineered differences and rolling statistics.

Workshop 2 asks whether a deep model can exploit the temporal structure of
the episodes directly. The central question is: \emph{which architecture,
training procedure and regularisation strategy turn the raw sequence of states
into a better failure predictor?}

The contributions of this report are: (i) an architecture selection justified
by the domain (recurrent networks over past-only windows); (ii) a
leakage-free sequence pipeline that reuses the Workshop 1 episode split and
scales with training statistics only; (iii) a training procedure adapted to a
17.2:1 imbalance (position-weighted loss) with modern regularisation; and
(iv) a reproducible comparison against the Workshop 1 baselines on exactly the
same test rows.

Section~II describes the architecture and its implementation;
Section~III details the data pipeline and the training protocol;
Section~IV reports the results and the comparison with Workshop 1; and
Section~V concludes.

%% ======================================================================
\section{Architecture Selection and Model Description}

\subsection{Why a recurrent network?}
The data are sequences: each episode is a time series of
$(x, \dot x, \theta, \dot\theta, a)$ and the target of a row is defined by the
\emph{future} evolution of the system. Among the architectures listed in the
workshop brief:
\begin{itemize}
\item \textbf{CNNs} are designed for images or spatial data; a 1-D convolution
      would read a fixed-length local patch but has no notion of the evolving
      hidden state that a control signal produces.
\item \textbf{Transformers} excel at long-range dependencies in large corpora,
      but with only 400 episodes their data hunger and positional machinery
      are overkill (and a source of overfitting).
\item \textbf{Value-based networks (DQN)} estimate action values for a
      reinforcement-learning agent; here we solve a \emph{supervised}
      prediction problem, not control.
\item \textbf{Recurrent networks (LSTM/GRU)} maintain a hidden state across
      time steps, which is precisely the inductive bias needed to integrate
      the recent history of the pole and detect the drift towards failure.
\end{itemize}
No pretrained model exists for this domain, so (as the brief indicates) the
network is trained from scratch with careful initialisation (Xavier for the
input gates, orthogonal for the recurrent weights, forget-gate bias $=1$) and
a learning-rate schedule.

\subsection{Model}
Fig.~\ref{fig:model} summarises the architecture implemented in
\texttt{src/w2/model.py}:

\begin{itemize}
\item \textbf{Input:} a window $(B, T, F)$ with $T=16$ past steps and
      $F=5$ features per step ($x$, $\dot x$, $\theta$, $\dot\theta$,
      action), z-score standardised with training statistics.
\item \textbf{Encoder:} \texttt{pack\_padded\_sequence} masks the padding of
      short episodes, so the final hidden state always corresponds to the
      last \emph{real} step; a stacked LSTM with 2 layers $\times$ 64 hidden
      units, dropout 0.2 between layers, unidirectional (the window already
      contains only past information with respect to the predicted row).
\item \textbf{Head:} Linear$(64 \rightarrow 32)$ -- ReLU -- Dropout$(0.2)$ --
      Linear$(32 \rightarrow 1)$ producing one logit per window.
\end{itemize}

\begin{figure}[t]
\centering
\figph{architecture diagram placeholder (input window, LSTM, MLP head)}{fig2-architecture}
\caption{Architecture of the deep failure predictor: a past-only window of 16
steps is encoded by a stacked LSTM; the final hidden state feeds a small MLP
head that outputs one logit.}
\label{fig:model}
\end{figure}

The model has \textbf{53{,}569 trainable parameters} (LSTM
$4H(F+H)+4H^2$ per layer $\approx$ 49.8k for two layers, plus a 2.1k head).
Alternatives considered: a GRU (fewer parameters, 40{,}705 -- reported in the
ablation), a bidirectional encoder (139{,}841 parameters; the reverse
direction only re-reads past steps, so it adds capacity but no new
information) and a larger 128-unit model. The forward-pass smoke test
(\texttt{python src/w2/model.py}) verifies the output shape and the parameter
count.

%% ======================================================================
\section{Data Pipeline and Training Protocol}

\subsection{Past-only sliding windows}
Every raw row becomes one training sample described by the last
$\min(t{+}1, 16)$ steps of the \emph{same} episode, including the current one
(\texttt{src/w2/sequences.py}). Three properties preserve the Workshop 1
rigour:
\begin{enumerate}
\item \textbf{No future information:} the window ends at the predicted row.
\item \textbf{Same test rows:} one sample per raw row, so the deep model is
      evaluated on exactly the 5{,}330 test rows (290 positives) used by the
      baselines.
\item \textbf{Leakage-free scaling:} the z-score mean and standard deviation
      are computed on the training partition only and stored inside the
      checkpoint.
\end{enumerate}
Short episodes simply produce shorter windows; the packed representation
masks the right-side padding (which is filled with zeros, i.e.\ the training
mean after standardisation). The episode-level split of
\texttt{data/splits.json} (279/61/60 episodes, seed 42) is reused verbatim.

\subsection{Handling class imbalance}
With 5.5\% positives, a model that always predicts ``safe'' reaches 94.5\%
accuracy. We therefore:
\begin{itemize}
\item weight the positive class in the loss with
      $\text{pos\_weight} = N_{neg}/N_{pos} = 16.96$ (computed on the
      training partition only), which is equivalent to class weighting in the
      objective function;
\item select the decision threshold on the \emph{validation} set by maximising
      F1, instead of assuming 0.5;
\item report PR-AUC and recall as the primary metrics, not accuracy.
\end{itemize}

\subsection{Optimisation, regularisation and monitoring}
\begin{itemize}
\item \textbf{Loss:} BCE-with-logits with the position weight above.
\item \textbf{Optimiser:} AdamW (learning rate $10^{-3}$, weight decay
      $10^{-4}$), batch size 256.
\item \textbf{Schedule:} cosine annealing from $10^{-3}$ down to
      $10^{-5}$ over the epochs.
\item \textbf{Regularisation:} dropout 0.2 (RNN layers and head), weight
      decay, gradient clipping at norm 1.0 (RNNs are prone to exploding
      gradients) and \textbf{early stopping} (patience 8 on validation F1).
\item \textbf{Monitoring:} every epoch is logged to TensorBoard
      (\texttt{runs/workshop2/tb/}) \emph{and} to a CSV (\texttt{runs/workshop2/train\_log\_*.csv})
      with train/val loss, F1, precision, recall, ROC-AUC, PR-AUC, accuracy,
      the selected threshold and the learning rate, so the learning curves of
      Fig.~\ref{fig:curves} can be regenerated without TensorBoard.
\item \textbf{Checkpointing:} the epoch with the best validation F1 is saved
      to \texttt{checkpoints/workshop2/model\_best.pt} together with the scaler, the
      window configuration and the threshold, making the artifact
      self-contained for \texttt{src/w2/evaluate\_lstm.py} and
      \texttt{src/w2/test\_model.py}.
\end{itemize}
Seeds are fixed to 42 (Python, NumPy and \texttt{torch}), so the runs are
reproducible.

%% ======================================================================
\section{Results and Discussion}

\subsection{Training dynamics}
Fig.~\ref{fig:curves} shows the learning curves of the selected run
(LSTM, window 16, raw features). The training loss decreases smoothly and the
validation loss follows it until epoch $\approx$15, after which the validation
loss plateaus while the training loss keeps decreasing -- the classic onset of
overfitting that triggers early stopping. The validation F1 climbs from 0.76
(epoch 1) to \textbf{0.881} at epoch 15 (best epoch), with PR-AUC 0.957; the
selected threshold stabilises around 0.9.

\begin{figure*}[t]
\centering
\figph[0.95\textwidth]{learning curves: loss, validation metrics and threshold per epoch}{fig2-curves}
\caption{Training and validation loss, validation metrics (F1, PR-AUC,
ROC-AUC) and the selected decision threshold per epoch for the best run
(LSTM, window 16, raw features, seed 42). The dashed line marks the best
epoch (15) selected by validation F1.}
\label{fig:curves}
\end{figure*}

\subsection{Test-set evaluation}
Table~\ref{tab:comparison} compares the Workshop 1 baselines with the deep
model on the same 60 test episodes at the validation-selected threshold
(confusion matrix in Fig.~\ref{fig:cm2}, ROC/PR curves in
Fig.~\ref{fig:rocpr}).

\begin{table}[t]
\caption{WORKSHOP 1 BASELINES VS. WORKSHOP 2 DEEP MODEL (TEST SET, VAL THRESHOLD)}
\label{tab:comparison}
\centering
\begin{tabular}{|l|c|c|c|c|c|}
\hline
\textbf{Model} & \textbf{Prec.} & \textbf{Rec.} & \textbf{F1} & \textbf{ROC-AUC} & \textbf{PR-AUC} \\ \hline
Logistic Regression & 0.669 & 0.703 & 0.686 & 0.981 & 0.791 \\ \hline
Random Forest & 0.812 & 0.852 & 0.832 & 0.994 & 0.923 \\ \hline
Gradient Boosting & 0.834 & 0.886 & 0.860 & 0.996 & 0.935 \\ \hline
\textbf{Deep LSTM (W2)} & \textbf{0.812} & \textbf{0.924} & \textbf{0.865} & \textbf{0.997} & \textbf{0.956} \\ \hline
\end{tabular}
\end{table}

\begin{figure*}[t]
\centering
\figph[0.95\textwidth]{bar chart comparing W1 baselines and the W2 deep model}{fig2-comparison}
\caption{Precision, recall, F1, ROC-AUC and PR-AUC on the test set: the three
Workshop 1 baselines versus the Workshop 2 deep model (validation threshold).}
\label{fig:compare}
\end{figure*}

\begin{figure}[t]
\centering
\figph{confusion matrix of the deep model on the test set}{fig2-confusion}
\caption{Confusion matrix of the deep model on the test set at the
validation threshold (TN = 4{,}978, FP = 62, FN = 22, TP = 268).}
\label{fig:cm2}
\end{figure}

\begin{figure*}[t]
\centering
\figph[0.95\textwidth]{ROC and precision-recall curves of the deep model and the baselines}{fig2-roc-pr}
\caption{ROC and Precision--Recall curves on the test set: the deep model
(solid blue) against the three Workshop 1 baselines (dashed).}
\label{fig:rocpr}
\end{figure*}

The deep model is the best classifier on every ranking metric: it detects
\textbf{268 of the 290} imminent failures (recall 0.924, versus 257 for
Gradient Boosting) with 62 false alarms, and its PR-AUC of 0.956 exceeds the
0.935 of the best baseline. Its F1 (0.865) is only marginally above the 0.860
of Gradient Boosting, because precision drops slightly (0.812 vs.\ 0.834): the
gain in recall costs some false positives. At threshold 0.5 the model reaches
recall 0.972 with precision 0.689 (F1 0.807), again illustrating that the
threshold, not the network, controls the precision/recall trade-off.

\subsection{What did the model learn?}
The network uses \emph{only the five raw variables} -- no distances to the
failure limits, no rolling statistics -- yet it outperforms the baselines
that were fed the 16 engineered features. This is evidence that the LSTM
recovered, from the temporal context itself, representations similar to the
hand-crafted ones (the drift of $|\theta|$ and $|\dot\theta|$ towards the
12$^\circ$ limit and the approach speed towards $|x| = 2.4$ m are exactly the
kind of patterns a hidden state integrates over 16 steps). The remaining
errors concentrate on states a few steps \emph{before} the label window,
where the label boundary is arbitrary with respect to the dynamics -- the
same phenomenon hypothesised in Workshop 1.

\subsection{Hyperparameters and ablations}
Table~\ref{tab:ablation} summarises the configurations tried (each trained
with the same protocol, seed 42 and early stopping). Four recurrent
architectures were compared: the cell (LSTM vs.\ GRU), the history length
(window 16 vs.\ 32) and the capacity (64 vs.\ 128 hidden units).

\begin{table}[t]
\caption{ABLATION STUDY: VALIDATION AND TEST METRICS (SEED 42)}
\label{tab:ablation}
\centering
\begin{tabular}{|l|c|c|c|c|}
\hline
\textbf{Configuration} & \textbf{Params} & \textbf{Val F1} & \textbf{Test F1} & \textbf{Test PR-AUC} \\ \hline
LSTM, window 16, 64h (main) & 53{,}569 & 0.881 & 0.865 & 0.956 \\ \hline
GRU, window 16, 64h & 40{,}705 & 0.886 & 0.850 & 0.954 \\ \hline
LSTM, window 32, 64h & 53{,}569 & 0.877 & 0.870 & 0.956 \\ \hline
LSTM, window 16, 128h & 205{,}377 & 0.885 & 0.852 & 0.956 \\ \hline
W1 Gradient Boosting (reference) & --- & --- & 0.860 & 0.935 \\ \hline
\end{tabular}
\end{table}

The validation F1 of the four runs spans only 0.877--0.886 (range 0.009),
which is smaller than the standard error of an F1 estimate on our validation
partition ($\approx$0.015 with 300 positive samples): the validation set
cannot reliably distinguish these configurations. Consistently, the
validation and test rankings are \emph{inversely} ordered (the GRU wins on
validation with 0.886 yet scores the worst test F1, 0.850), so selecting the
final model by the validation maximum would amount to fitting noise, and
selecting by the test set would be leakage. We therefore report the
\textbf{pre-specified main configuration} (LSTM, window 16, 64 hidden units,
chosen before running the variants) as the final model and disclose every
run in Table~\ref{tab:ablation}. All four variants beat or match the best
Workshop 1 baseline on PR-AUC (0.954--0.956 vs.\ 0.935), and three of the
four also beat it on F1 (0.850--0.870 vs.\ 0.860), confirming that the gain
comes from the recurrent architecture rather than from a lucky
hyperparameter choice. Doubling the hidden units (205k parameters) does not
help, and the GRU trades 13k fewer parameters for a small F1 drop -- a useful
efficiency result. We did not run the 16 engineered features: the point of
the deep model is to learn its own temporal representations from the raw
states.

\subsection{Limitations}
As in Workshop 1, all numbers come from a single episode-level split with one
seed and 60 test episodes; the difference between F1 0.865 and 0.860 is within
the noise of such a small test set, so we claim \emph{comparability with a
consistent improvement on recall and PR-AUC}, not a decisive win. Grouped
cross-validation by episode and a multi-seed study would quantify the
variance. The threshold is re-selected on validation every epoch, which adds
a mild optimism to the validation score (the test score, however, is computed
once with the frozen threshold).

%% ======================================================================
\section{Conclusions}
We selected and implemented a recurrent architecture appropriate for the
CartPole failure-prediction domain: a stacked LSTM over past-only windows of
16 steps, trained from scratch with position-weighted BCE, AdamW, cosine
scheduling, dropout, weight decay, gradient clipping and early stopping, and
monitored with TensorBoard and CSV logs. On the same test rows used in
Workshop 1 the model reaches F1 = 0.865, PR-AUC = 0.956 and recall = 0.924,
improving every ranking metric over the best baseline (Gradient Boosting,
F1 = 0.860, PR-AUC = 0.935, recall = 0.886) while consuming only the raw
state variables. The best model is saved in
\texttt{checkpoints/workshop2/model\_best.pt} and can be reloaded independently of the
training code. Future work includes grouped cross-validation, multi-seed
evaluation, and fine-tuning strategies such as unfreezing layers with a lower
learning rate if a larger corpus becomes available.

\section*{Reproducibility}
Code, split file and instructions are available at
\url{https://github.com/Tonyaa11/Cartpole}. To reproduce the results:
\texttt{python src/w2/train\_lstm.py}, then \texttt{python src/w2/evaluate\_lstm.py}, then
\texttt{python src/w2/make\_figures\_w2.py} (Python 3.11.2, PyTorch 2.6.0 CPU,
scikit-learn 1.6.1 as in Workshop~1, seed 42). The per-epoch logs are in
\texttt{runs/workshop2/train\_log\_*.csv} and the comparison table in
\texttt{runs/workshop2/comparison\_*.csv}.

\begin{thebibliography}{00}
\bibitem{barto1983} A. G. Barto, R. S. Sutton, and C. W. Anderson, ``Neuronlike adaptive elements that can solve difficult learning control problems,'' \emph{IEEE Trans. Syst., Man, Cybern.}, vol. SMC-13, no. 5, pp. 834--846, 1983.
\bibitem{hochreiter1997} S. Hochreiter and J. Schmidhuber, ``Long short-term memory,'' \emph{Neural Comput.}, vol. 9, no. 8, pp. 1735--1780, 1997.
\bibitem{kingma2015} D. P. Kingma and J. Ba, ``Adam: A method for stochastic optimization,'' in \emph{Proc. ICLR}, 2015.
\bibitem{loshchilov2019} I. Loshchilov and F. Hutter, ``Decoupled weight decay regularization,'' in \emph{Proc. ICLR}, 2019.
\bibitem{paszke2019} A. Paszke et al., ``PyTorch: An imperative style, high-performance deep learning library,'' in \emph{NeurIPS}, 2019.
\bibitem{saito2015} T. Saito and M. Rehmsmeier, ``The precision-recall plot is more informative than the ROC plot when evaluating binary classifiers on imbalanced datasets,'' \emph{PLoS ONE}, vol. 10, no. 3, e0118432, 2015.
\end{thebibliography}

\end{document}
